% SPDX-License-Identifier: CC-BY-4.0
% Copyright (C) 2026 Fabian Januszewski, Abhijit (the Langlandscape team)
% Stage 05, WP S05-M7 (package abhijit/s05-math): which indefinite stage appears (D1), the two pairings (BN), context, the
% simplicial shape and the two-part zig-zag, the open points O2, O4, O6, O7, and the degree remark.
% Labels: ss:s05-simplicial, rem:s05-D1, rem:s05-BN, rem:s05-context, rem:s05-degree, def:s05-stages, conj:s05-O2, rem:s05-O7,
% rem:s05-O46.
\subsection{Simplicial shape, the zig-zag and open points}\label{ss:s05-simplicial}

\begin{remark}[which ``indefinite stage'' appears; the D1 seam]\label{rem:s05-D1}
$\HypX$ compares $\Xq(J_s)$ with $\Z[S_s]^0=\Xcomb_q(U_s)$ (\ref{if:s03-graph}), the \emph{graph} homology $H_1(|\GDR_q(U_s)|,\Z)$ of the
two-vertex DR graph; it is torsion-free, as is every character group. Wherever this stage uses a tree model at a \emph{second} prime $r$
(the $\GL_2$ form of the Ribet sequence, \S\ref{sss:s05-general}), it uses the amended \pathtt{if:s02-indef}: over $\Z$ the indefinite stage is
$\cX^{(r)}(U):=H_1(\BG(U,r),\Z)=H_1(\Gamma_U,\Z)/\text{torsion}$, \emph{not} $H_1(\Gamma_U,\Z)$, whose torsion (of exponent dividing $24$) has no
counterpart in a character group. Over $\Zp$, $p\ge5$, the two agree (stage 02, D1). No statement of this stage uses
$H_1(\Gamma_U,\Z)_{\rm tors}$. \stproved{} (bookkeeping; the torsion statement is stage 02's).
\end{remark}

\begin{remark}[$\eta^B$ and (X2): two different pairings]\label{rem:s05-BN}
Stage 04 (D4) computes $\eta^B_F$ with Hsieh's pairing $B_N$: a fixed-character piece, an Atkin--Lehner operator at $Mp^\alpha$, a norm
twist, a sum over $X_0(Np^\alpha)$ (level $\Uzero{Mp^\alpha}$, Hsieh's $N=M$) with weights $(\#\Gamma_{Np^\alpha,x})^{-1}$ on the function side
\lit{hsieh\_2017\_hida-families-and-p-adic-triple-product-l-functions\_arxiv}{(4.1), p.~29} (stage 04, item (iv) after Def.~BN), and a sign $(-1)^k$.
$\langle,\rangle_w$ is a sum over $S_s$ (level $\Uone{Mp^s}$) with $w_x=1$; unfolding from $\Uzero{Mp^s}$ to $S_s$ introduces the constant $\varphi(Mp^s)/2$
(Remark~\ref{rem:s05-c2open}(vi)), a $p$-adic non-unit for $s\ge2$, so the two normalisations differ by more than a unit.
(X2) uses $\langle,\rangle_w$ without Atkin--Lehner. Since $\langle x,w\,y\rangle_w$ is $\theta$-equivariant exactly when $\langle,\rangle_w$ is
$\theta^\Ros$-equivariant (Prop.~\ref{prop:s05-AL}), the table row ``$\eta^B_F=$ monodromy self-pairing via (X2)'' (\S\ref{ss:s05-lost}) means:
$\eta^B_F$ should be generated by $[g,w_{Mp^s}g]_s$ (the form of Prop.~\ref{prop:s05-Xlambda}(c), i.e.\ $u_s(g,w_{Mp^s}\cdot)$ with the $U_p^{-s}$-twist, equal to $u_s(g,w_{Mp^s}g)$ up to a $p$-adic unit on an eigenline; $[\,,\,]_s$ has no norm-twist factor, which sits in the operator $A_s$ of (c2)(i),(v)) for a generator $g$ of the $F$-eigenline of $\Xq(J_s)_\m$, up to a constant: the constant $\varphi(Mp^s)/2$ of Remark~\ref{rem:s05-c2open}(vi) is \emph{open} (not claimed to be a unit; the sign
$(-1)^k$ is a unit). \sttoprove{} (bookkeeping against stage 04's normalisation of $B_N$).
\provisional{S05-M10: this reading depends on (c2) of Prop.~\ref{prop:s05-Xlambda} (matching $u_s(\cdot,w_{Mp^s}\cdot)$ with $B_N$), which is \stconj{} and
\stopen{} (Remark~\ref{rem:s05-c2open}); until (c2) is settled the identification of $\eta^B_F$ with a monodromy self-pairing is not claimed.}
\end{remark}

\begin{remark}[context]\label{rem:s05-context}
That an integral comparison must produce a defect is the point of
\lit{ichino-prasanna\_2016\_periods-of-quaternionic-shimura-varieties-i\_arxiv}{Abstract, p.~1; Introduction, items (i)--(ii), p.~2: the period ratio is
an explicit product of Tamagawa numbers, related to level-lowering congruences}; a comparable integral
(torsion) Jacquet--Langlands in a different setting is
\lit{calegari-venkatesh\_2012\_torsion-jacquet-langlands-correspondence\_arxiv}{\S1.1.1, pp.~2--3, and Theorem~A, p.~3}, which also compares orders of finite homology groups rather than producing a map. \stcited.
\end{remark}

\begin{remark}[no Betti-canonical replacement]\label{rem:s05-degree}
After inverting $q$, the cuspidal cohomology of $\GL_2$ sits in degree $2$ (\ref{if:s03-amalgam}: the $q$-new part of
$H^*(\Gamma_0(M)[\frac1q],V)$ is in $H^2$) and that of the definite $S$-arithmetic group in degree $1$ (\texttt{if:s02-coh},
Prop.~\ref{prop:s02-coh}); for $B_q$ itself ($q\mid D$) there is no tree and it sits in degree $0$ (Prop.~\ref{prop:s05-nochain}(a)). A
degree-preserving Hecke-equivariant chain map inducing $\HypX$ on cuspidal parts therefore cannot exist between the $S$-arithmetic
complexes of $\GL_2(\Z[\frac1q])$-type (cuspidal part in degree $2$) and of $B_q$ (degree $0$); any realisation must shift degree (e.g.\ through a boundary or Gysin-type map) or go through a third object. \stsketch{}
(degree count only).
\end{remark}

\begin{definition}[the four stages as simplicial orbifolds]\label{def:s05-stages}
\begin{enumerate}[label=(\roman*)]
\item Definite level $U$: the $0$-dimensional orbifold (groupoid) $\coprod_{[x]\in\XB{U}}B\Gamma_x$; its $H^0$ with coefficients $V$ is
  $\Brandt{U}{V}$.
\item Definite $S$-arithmetic stage: the graph of finite groups $\Gamma_U\quot\cT_q$ (\texttt{if:s02-quot}, Prop.~\ref{prop:s02-quot}) with
  structure maps $s_0,s_1$; a $1$-dimensional orbifold model of $B\Gamma_U$ after barycentric subdivision, in the sense of orbifolds as
  étale groupoids \lit{moerdijk-pronk\_1997\_orbifolds-sheaves-and-groupoids}{\S4, Thm.~4.1, p.~11: orbifolds correspond to effective
  étale groupoids with proper $(s,t)$} ; for the simplicial model of an orbifold and its cohomology with local coefficients:
  \lit{moerdijk-pronk\_1999\_simplicial-cohomology-of-orbifolds}{Theorem in the introduction, p.~269; Thm.~2.1.1, p.~278}. Over $\Z$ its $H_1$ is taken as in
  Remark~\ref{rem:s05-D1}.
\item $\GL_2$ arithmetic stage: $\Gamma_s\backslash\Hb$ with a finite triangulation (as in HACS).
\item $\GL_2$ $S$-arithmetic stage: the graph of spaces $[Y_0(M)\sqcup Y_0(M)\leftleftarrows Y_0(Mq)]$ (\ref{if:s03-amalgam}).
\end{enumerate}
\end{definition}
The zig-zag
\[\underbrace{\XB{U}\rightleftarrows\XB{\Uzero{q}U}}_{\text{0-dim}}\ \leadsto\ \underbrace{\Gamma_U\quot\cT_q}_{\text{1-dim}}
\qquad\qquad \underbrace{\XBq{U_s}}_{\text{edges of }\GDR_q}\ \overset{\HypX}{\longleftrightarrow}\ \underbrace{\text{dual graph of }X_{s,\bar\F_q}}_{\GL_2}\]
has two independent parts: the first (for $B$ with $q\nmid D$) consists of explicit simplicial maps; the second (for $B_q$) is $\HypX$,
a statement that the DR graph $\GDR_q(U_s)$ of stage 03 (\ref{if:s03-graph}) \emph{is} the dual graph of the special fibre, with
thicknesses, compatibly with Hecke and Galois. (The seed's single displayed zig-zag concatenates the two, which is misleading: the
left part is at a prime $q\nmid D$, the right part for $D=q$.) At $\Gamma_0$-level the right-hand identification of the dual graph is
Ribet's description (two components, edges $=$ supersingular points) \lit{ribet\_1990\_modular-representations-of-gal-qbar-q-arising-from-modular-forms}{p.~439}.

\begin{conjecture}[O2, a simplicial realisation; open]\label{conj:s05-O2}
There is a Hecke-equivariant map of finite simplicial (orbi-)sets, defined over $\Z[\frac1{6}]$ at most, from $\GDR_q(U_s)$ to a simplicial
model of the special fibre of $X_s$ at $q$ (or to a $q$-adic tubular neighbourhood), inducing $\varphi_s^{-1}$ on $H_1$. \stconj{};
equivalently (by Remark~\ref{rem:s05-degree}) a degree-shifting map between the two $S$-arithmetic complexes. The most promising concrete
form is the Deuring map on supersingular points, which exists only after reduction mod $q$; we expect no characteristic-zero Betti
version. \stopen.
\end{conjecture}

\begin{remark}[O7: relation to Chenevier]\label{rem:s05-O7}
Chenevier's $\JL_p$ is a closed immersion of eigencurves
\lit{chenevier\_2005\_une-correspondance-de-jacquet-langlands-p-adique}{\S4.6, Thm.~3, p.~188}. On the ordinary locus, the maps of
Prop.~\ref{prop:s05-Xlambda}(d) induce on Hecke algebras the surjection $\T^{\ord}(qM)\twoheadrightarrow\T^{\ord}_{B_q}(M)$ of
\pathtt{if:s01-targets}, so their ``support'' is the ordinary part of the image of $\JL_p$. \stsketch. Beyond the ordinary locus nothing
is claimed (O7 stays open): there is no finite-slope analogue of $e$ killing $\Phicomb$.
\end{remark}

\begin{remark}[O4, O6 made precise]\label{rem:s05-O46}
O4: a Shimura-curve $\HypX$ for discriminant $d$ and $r\nmid d$ is Hypothesis~\ref{hyp:s05-X} with $J_s$ replaced by the Jacobian of the
Shimura curve of discriminant $d$ and level $\Gamma_1(Mp^s)\cap\Gamma_0(r)$, $B=B_{dr}$; at $\Gamma_0$-level it is
\lit{kohel\_2001\_hecke-module-structure-of-quaternions\_preprint}{Thm.~4.3, p.~11} (Buzzard). O6: at an Eisenstein $\m$, (X2) carries
information that $\HypXiso$ does not; its weight-$2$, level-$1$ shadow is the Mazur check $\Phicomb(\GDR_q(1))\cong\Z/\mathrm{num}\frac{q-1}{12}$
(\ref{if:s03-graph}) versus $\Phi_q(J_0(q))$ (\lit{mazur\_1977\_modular-curves-and-the-eisenstein-ideal}{Appendix, Thm.~(A.1)(b),
p.~173}); this is a check of (X2) at $\Gamma_0$-level, not a proof.
\end{remark}
